\subsection*{Inverse supervision and evidence localization}
We initialized two writers from the final forward participant checkpoint. Each of six epochs included all 19,426 ordinary training questions and two presentations of each of 2,508 human inverse questions. The training articles, passages and verbatim questions were disjoint from the 387-question development cohort. We also presented 652 within-passage competing pairs per epoch and generated closed rotations with two or three distinct participants. Constructed aliases comprised one to three printed or translated glyph cells. Names were absent from the question. The two selected bronze test images and the four further historical state glyphs were excluded from this alias pool. In the randomized arm the extra human question pairs exchanged targets with probability one half, and each synthetic scene received a random cyclic displacement of its targets. The two arms retained identical single-question examples, role and inference rehearsal, final checkpoint rule, and 2,430 optimizer steps. The three per-task masked negative log likelihoods for natural answering, inverse scenes and rehearsal had weights 1, 0.8 and 0.8. Both models started each response empty and freely selected source images, writing-memory images or stop.

For the four-participant test, four visible names and four objects formed a passage. The question named one object, and all four rotations of the names were tested for each object. We scored exact stopped image sequences and separately counted objects recovered in all rotations. The four images represent Jin and Cai in bronze script, Yan in oracle script and Wei in the Liushutong collection. Properties are invented for evaluation. For the DRCD recount we rendered each accepted human answer found in its source passage and compared the resulting images directly with the archived output. The failure categories for first glyph and contiguous source span overlap.

To vary the amount of visible evidence at fixed model weights, we split each pictured passage at glyphs for terminal punctuation. The oracle choice was the shortest sentence enclosing a reference. The negative control was a sentence without a reference that most closely matched its length. An inference-time selector chose the sentence maximizing $|Q\cap S|/|Q\cup S|$, where $Q$ and $S$ contain the distinct glyph images in the question and candidate sentence after removing terminal punctuation. Ties were resolved in reading order. All conditions retained the full question, decoding policy and original human reference set. Differences use 10,000 paired resamples of complete passages. The oracle condition measures what this writer can recover when a reference helps select the visible sentence.
